\documentclass[11pt]{article}
\usepackage[T1]{fontenc}
\usepackage[utf8]{inputenc}
\usepackage{lmodern}
\usepackage[margin=0.78in]{geometry}
\usepackage{amsmath,amssymb,amsthm}
\usepackage{booktabs}
\usepackage{microtype}
\usepackage{array}
\usepackage{titlesec}
\titlespacing*{\section}{0pt}{4pt}{1pt}
\titleformat*{\section}{\large\bfseries}
\setlength{\parskip}{1pt}
\setlength{\parindent}{0pt}
\allowdisplaybreaks
\newtheorem{proposition}{Proposition}
\newtheorem{corollary}{Corollary}

\title{\vspace{-3.5em}The Eigenstructure of the Black--Scholes Second-Order Greek Matrix\\[2pt]
\large A note on $(S,\sigma)$ curvature, its signature, and the at-the-money eigenframe}
\author{}
\date{2026-07-24}

\begin{document}
\maketitle
\vspace{-3.2em}

\begin{abstract}
\noindent The Hessian of the Black--Scholes price in spot and implied volatility has
signature governed by $\operatorname{sign} d_2$ alone; at the money it factors into a
scalar times a matrix of determinant $-\tfrac12$, so its eigenframe is fixed and vanna
acts only as a rotation by half the total variance. Eigenvectors are coordinate-dependent;
the inertia is not.
\end{abstract}

\vspace{-0.4em}
\section{Setup}
Work in the Black--Scholes frame with $r=q=0$, so $F=S$; fix strike $K$, maturity $T$, implied
volatility $\sigma$, and write $m=\ln(K/S)$. Then
\[
d_{1,2}=-\frac{m}{\sigma\sqrt T}\pm\frac{\sigma\sqrt T}{2},\qquad
d_1-d_2=\sigma\sqrt T,\qquad d_1 d_2=\tfrac14\sigma^2T-\tfrac{m^2}{\sigma^2T},
\]
and $\varphi=\varphi(d_1)$ is the standard normal density at $d_1$. The second-order sensitivities
in $(S,\sigma)$ assemble into the symmetric Hessian
\[
H=\begin{pmatrix}\Gamma&\text{Vanna}\\[1pt]\text{Vanna}&\text{Volga}\end{pmatrix},\qquad
\Gamma=\frac{\varphi}{S\sigma\sqrt T},\quad
\text{Vanna}=-\frac{\varphi\,d_2}{\sigma},\quad
\text{Volga}=\frac{S\sqrt T\,\varphi\,d_1 d_2}{\sigma}.
\]
$H$ is the curvature of the price as a function of the two state variables a desk marks against;
its eigenstructure is the local principal-axis decomposition of that curvature.

\section{Signature}
\begin{proposition}\label{prop:sig}
$\det H=\dfrac{\varphi^2\sqrt T}{\sigma}\,d_2$, hence $\operatorname{sign}\det H=\operatorname{sign} d_2$.
\end{proposition}
\begin{proof}
Directly,
\[
\det H=\Gamma\,\text{Volga}-\text{Vanna}^2
=\frac{\varphi^2 d_1 d_2}{\sigma^2}-\frac{\varphi^2 d_2^2}{\sigma^2}
=\frac{\varphi^2}{\sigma^2}\,d_2\,(d_1-d_2)
=\frac{\varphi^2}{\sigma^2}\,d_2\,\sigma\sqrt T
=\frac{\varphi^2\sqrt T}{\sigma}\,d_2. \qedhere
\]
\end{proof}
Since $\Gamma>0$ always, the signature is read off $d_2$: $H$ is positive definite when $d_2>0$,
rank one when $d_2=0$, and a saddle when $d_2<0$.
\begin{corollary}
$d_2>0\iff m<\tfrac12\sigma^2T\iff S>K\,e^{\sigma^2T/2}$. Curvature in $(S,\sigma)$ is a genuine
minimum only sufficiently far in the money; at spot--ATM ($m=0$) one has $d_2=-\tfrac12\sigma\sqrt T<0$,
so the ATM point lies strictly in the saddle region --- $\Gamma>0$ and $\text{Volga}>0$ coexist with a
negative eigenvalue produced by vanna.
\end{corollary}

\section{The at-the-money factorization}
Set $m=0$ (so $K=S$): then $d_1=-d_2=\tfrac12\sigma\sqrt T$ and $d_1d_2=-\tfrac14\sigma^2T$. Writing
$\beta:=S\sigma T$,
\[
H=\varphi\sqrt T\,M,\qquad
M=\begin{pmatrix}\dfrac1\beta&\dfrac12\\[6pt]\dfrac12&-\dfrac\beta4\end{pmatrix},\qquad
\det M=-\frac1\beta\cdot\frac\beta4-\frac14=-\frac12 .
\]
\begin{proposition}\label{prop:eig}
$\det M=-\tfrac12$ for every $\beta>0$. With $t:=\operatorname{tr}M=\dfrac1\beta-\dfrac\beta4$, the
eigenvalues and axes of $M$ are
\[
\lambda_\pm=\frac{t\pm\sqrt{t^2+2}}{2},\qquad
\lambda_+\lambda_-=-\tfrac12,\qquad
v_\pm\propto\Bigl(1,\;2\bigl(\lambda_\pm-\tfrac1\beta\bigr)\Bigr),\qquad
\tan2\theta=\frac{4\beta}{4+\beta^2},
\]
where $\theta$ is the rotation of the eigenframe from the $(S,\sigma)$ axes.
\end{proposition}
\begin{proof}
$\lambda_\pm$ solve $\lambda^2-t\lambda+\det M=0$ with $\det M=-\tfrac12$, so the discriminant is
$t^2-4\det M=t^2+2$ and $\lambda_+\lambda_-=\det M=-\tfrac12$. For the axis, $(M-\lambda I)v=0$ gives
$v\propto(M_{12},\,\lambda-M_{11})=(\tfrac12,\lambda-\tfrac1\beta)$, i.e. $(1,2(\lambda-\tfrac1\beta))$;
and $\tan2\theta=2M_{12}/(M_{11}-M_{22})=4\beta/(4+\beta^2)$ for the symmetric $2\times2$ case.
\end{proof}
The determinant invariant $\det M=-\tfrac12$ means the two axes carry eigenvalues of strictly opposite
sign for every $\beta$: the ATM curvature is a saddle whose principal-axis product is fixed independent
of spot, volatility, and maturity.

\section{Coordinate dependence and the natural (log) normalization}
$H$ mixes units ($\Gamma$ per $S^2$, Volga per $\sigma^2$, Vanna per $S\sigma$), so its eigenvectors are
\emph{not} invariant: a rescaling by $D=\operatorname{diag}(a,b)$ sends $H\mapsto D^{-1}HD^{-1}$, changing
both axes and eigenvalues. By Sylvester's law of inertia only the signature (Proposition~\ref{prop:sig})
survives. The natural choice measures curvature against \emph{relative} moves $dS/S$, $d\sigma/\sigma$ ---
the congruence $\tilde H=DHD$, $D=\operatorname{diag}(S,\sigma)$, so $\tilde H_{11}=S^2\Gamma$,
$\tilde H_{12}=S\sigma\,\text{Vanna}$, $\tilde H_{22}=\sigma^2\,\text{Volga}$.
\begin{proposition}\label{prop:tilde}
At the money, with $w:=\sigma^2T$ the total variance,
\[
\tilde H=S\sigma\varphi\sqrt T\;M_w,\qquad
M_w=\begin{pmatrix}\dfrac1w&\dfrac12\\[6pt]\dfrac12&-\dfrac w4\end{pmatrix},\qquad
\det M_w=-\tfrac12,\qquad
\tan2\theta=\frac{4w}{4+w^2}.
\]
The prefactor is exactly $S\sigma\varphi\sqrt T$ --- \emph{not} $\varphi\sqrt T$; the factor $S\sigma$
is the Jacobian of the log rescaling and must be carried.
\end{proposition}
\begin{proof}
At $m=0$: $\tilde H_{12}=S\sigma\cdot(-\varphi d_2/\sigma)=-S\varphi d_2=\tfrac12 S\sigma\varphi\sqrt T$,
fixing the prefactor $P=S\sigma\varphi\sqrt T$. Then $\tilde H_{11}=S^2\Gamma=S\varphi/(\sigma\sqrt T)=P/w$
and $\tilde H_{22}=\sigma^2\,\text{Volga}=-S\varphi\sigma^3T^{3/2}/4=-Pw/4$, so $\tilde H=P\,M_w$; the
angle formula follows from Proposition~\ref{prop:eig} with $\beta\to w$.
\end{proof}
So the log-normalized ATM curvature has the identical $M$-form with $\beta$ replaced by the total
variance $w=\sigma^2T$. For a small total variance $w\to0$, using $\det M_w=\lambda_+\lambda_-=-\tfrac12$:
\[
\lambda_+=\frac1w+\frac w4+O(w^3)\ \approx\ \frac1w\ \ (\text{gamma axis}),\qquad
\lambda_-=-\frac1{2\lambda_+}\ \approx\ -\frac w2\ \ (\text{volga axis}),\qquad
\theta\approx\frac w2 .
\]
The volga \emph{entry} is $-w/4$, but vanna's off-diagonal contributes an equal $-w/4$ to the eigenvalue,
so the volga-axis eigenvalue is $\lambda_-\approx-w/2$, twice the entry --- the exact invariant
$\lambda_+\lambda_-=-\tfrac12$ forces it. Interpretation: at the money the gamma and volga directions
nearly decouple ($\lambda_+\lambda_-$ fixed, axes nearly the coordinate axes), and vanna's entire effect
is to rotate the eigenframe by $\theta\approx w/2$, half the total variance.

\section{Numerical check}
Take $S=K=100$, $\sigma=0.2$, $T=1$, so $\beta=S\sigma T=20$ and $\varphi\sqrt T=0.396953$; closed form
and \texttt{numpy.linalg.eigh} agree.
\begin{center}\small
\begin{tabular}{lll}
\toprule
Quantity & Closed form & Numeric \\
\midrule
$H_{11}=\Gamma$ & $\varphi\sqrt T/\beta=0.019848$ & $0.019848$\\
$H_{12}=\text{Vanna}$ & $\varphi\sqrt T\cdot\tfrac12=0.198476$ & $0.198476$\\
$H_{22}=\text{Volga}$ & $-\varphi\sqrt T\cdot\tfrac\beta4=-1.984763$ & $-1.984763$\\
$\det H$ & $(\varphi^2\sqrt T/\sigma)\,d_2=-0.078786$ & $-0.078786$\\
$\lambda_-,\ \lambda_+$ & $\varphi\sqrt T\cdot\{\tfrac{t\pm\sqrt{t^2+2}}2\}$ & $-2.004225,\ \ 0.039310$\\
\bottomrule
\end{tabular}
\end{center}
Here $\det H<0$ confirms the ATM saddle (Proposition~\ref{prop:sig}, $d_2=-0.1<0$), and
$\lambda_-\lambda_+=(\varphi\sqrt T)^2\det M=-\tfrac12(\varphi\sqrt T)^2$ recovers the invariant.
\end{document}
